\documentclass{article}
\usepackage{amsmath, amssymb, amsthm}
\usepackage{hyperref}
\usepackage{geometry}
\geometry{margin=1in}

\title{Erd\H{o}s Problem \#855}
\date{Last updated: 2026-03-14}
\author{Source: erdosproblems.com}

\begin{document}
\maketitle

\noindent\textbf{Status:} OPEN \quad \textbf{No prize}

\noindent\textbf{Tags:} number theory, primes

\medskip
\noindent\textbf{Note:} second Hardy-Littlewood conjecture

\medskip

\noindent\textbf{Problem Statement:}

If $\pi(x)$ counts the number of primes in $[1,x]$ then is it true that (for large $x$ and $y$)\[\pi(x+y) \leq \pi(x)+\pi(y)?\]




Commonly known as the second Hardy-Littlewood conjecture . In \cite{Er85c} Erd\H{o}s describes it as 'an old conjecture of mine which was probably already stated by Hardy and Littlewood'. This is probably false, since Hensley and Richards \cite{HeRi73} have shown that this is false assuming the Hardy-Littlewood prime tuples conjecture . Indeed, assuming this conjecture, they prove that for every large $y$ there are infinitely many $x$ such that\[\pi(x+y)&gt;\pi(x)+\pi(y)+(\log 2-o(1))\frac{y}{(\log y)^2},\]where the $o(1)$ term tends to $0$ as $y\to \infty$. Erd\H{o}s \cite{Er85c} reports Straus as remarking that the 'correct way' of stating this conjecture would have been\[\pi(x+y) \leq \pi(x)+2\pi(y/2).\]Clark and Jarvis \cite{ClJa01} have shown this is also incompatible with the prime tuples conjecture. In \cite{Er85c} Erd\H{o}s conjectures the weaker result (which in particular follows from the conjecture of Straus) that\[\pi(x+y) \leq \pi(x)+\pi(y)+O\left(\frac{y}{(\log y)^2}\right),\]which the Hensley and Richards result shows (conditionally) would be best possible. Richards conjectured that this is false. Erd\H{o}s and Richards further conjectured that the original inequality is true almost always - that is, the set of $x$ such that $\pi(x+y)\leq \pi(x)+\pi(y)$ for all $y&lt;x$ has density $1$. They could only prove that this set has positive lower density. They also conjectured that for every $x$ the inequality $\pi(x+y)\leq \pi(x)+\pi(y)$ is true provided $y \gg (\log x)^C$ for some large constant $C&gt;0$. In \cite{Er80} Erd\H{o}s conjectures that\[\pi(x+y)\leq \pi(x)+O\left(\frac{y}{\log x}\right)\]for every $x\geq y \gg (\log x)^C$ for some large constant $C&gt;0$. Hardy and Littlewood proved\[\pi(x+y) \leq \pi(x)+O(\pi(y)).\]The best known in this direction is a result of Montgomery and Vaughan \cite{MoVa73}, which shows\[\pi(x+y) \leq \pi(x)+2\frac{y}{\log y}.\]This is discussed in problem A9 of Guy's collection \cite{Gu04}.




References

[ClJa01] Clark, David A. and Jarvis, Norman C., Dense admissible sequences . Math. Comp. (2001), 1713--1718.

[Er80] Erd\H{o}s, Paul, A survey of problems in combinatorial number theory . Ann. Discrete Math. (1980), 89-115.

[Er85c] Erd\H{o}s, P., On some of my problems in number theory I would most like to see solved . Number theory (Ootacamund, 1984) (1985), 74-84.

[Gu04] Guy, Richard K., Unsolved problems in number theory . (2004), xviii+437.

[HeRi73] Hensley, Douglas and Richards, Ian, On the incompatibility of two conjectures concerning primes . (1973), 123--127.

[MoVa73] Montgomery, H. L. and Vaughan, R. C., The large sieve . Mathematika (1973), 119--134.

\noindent\textbf{Related OEIS sequences:} A023193


\bigskip
\noindent\small{Source: \url{https://www.erdosproblems.com/855}}
\end{document}
